% Final §1 -- Motivation and business case (plan §7.1: 1.25 pages with the title block).
% Market-sizing inputs: docs/market-sizing.md §2, §6 (keys in references.bib, notes in sources.md).
\section{Motivation and business case}\label{sec:motivation}

\paragraph{Problem.}
Many workshops on campus and in continuing education and training (CET) run the same content in
several sessions so that people with different timetables can attend. Registration is usually one
form per session plus a spreadsheet. A participant who could attend two sessions registers for
both to be safe; the organiser later has to find and release the duplicate seat by hand, and a seat
freed by a dropout stays empty unless someone e-mails the next person on a waiting list. Duplicate
holds block seats that other people want, and the reconciliation is manual work for the organiser.

\paragraph{SeatRelay.}
A participant registers \emph{once} per workshop with an ordered list of acceptable sessions and
joins \emph{one} queue; an account holds at most one seat per workshop. A free seat is offered to
the first queued person who accepts that session, and the offer is held for a time limit
(\code{holdSeconds}, default 900\,s). A seat released by an expiry, a decline, a cancellation or a
capacity increase is offered to the next person in the same decision (Fig.~\ref{fig:flow}). The
organiser manages 2--5 sessions and sees, per session, capacity, confirmed and held seats, the queue,
a backfill log with latencies, distinct participants served, attendance and a CSV export.
Participants see their status and queue position and receive in-app and e-mail notifications.

\begin{figure}[tb]
  \centering% Fig. 1 (plan §7.1): user flow of SeatRelay (participant states, backfill, organiser view).
% NOTE FOR THE TEAM: this is a redrawn stand-in. If you still have the source of Fig. 1 of the
% Preliminary report, use that figure instead (plan §7.1 says "Fig 1 flow (reuse)") and delete
% this file; only adapt its wording to the tick-first expiry and the polling protocol.
% Source of truth: docs/ARCHITECTURE.md §1 (overview), §9 (rules R1–R7).
% Needs % SeatRelay report: shared TikZ setup for the hand-drawn figures.
%
% Load once in the report preamble:   % SeatRelay report: shared TikZ setup for the hand-drawn figures.
%
% Load once in the report preamble:   \input{figures/tikz-setup}
% Then place a figure with
%   \begin{figure}[t]\centering\input{figures/fig-architecture}\caption{...}\label{fig:arch}\end{figure}
%
% The figure files contain only a tikzpicture (no figure environment, no caption), so captions
% and labels stay in the report text. All diagram text is set at a fixed 7 pt sans size, so the
% figures have the same physical size whatever the document's base font size is. Widths are
% checked against a single-column A4 layout with 2 cm margins (\textwidth = 17 cm): every
% figure is at most 0.95\textwidth wide.
%
% Colours: the categorical slots of the validated palette used by the result charts
% (report/tools/make_figures.py), as light tints for node fills. Identity never rests on colour
% alone: every node names its service in text.

\usepackage{tikz}
\usetikzlibrary{arrows.meta,positioning,calc,fit,backgrounds,shapes.geometric,decorations.pathreplacing}

\definecolor{srBlue}{HTML}{2A78D6}
\definecolor{srOrange}{HTML}{EB6834}
\definecolor{srGreen}{HTML}{1BAF7A}
\definecolor{srYellow}{HTML}{EDA100}
\definecolor{srPurple}{HTML}{6250D6}
\definecolor{srInk}{HTML}{1F1F1D}
\definecolor{srMuted}{HTML}{52514E}
\definecolor{srGrid}{HTML}{D8D7D2}

\newcommand{\srfont}{\sffamily\fontsize{7}{8.2}\selectfont}
\newcommand{\srfontb}{\sffamily\bfseries\fontsize{7}{8.2}\selectfont}

\tikzset{
  sr/base/.style={draw=srInk!75, line width=0.5pt, rounded corners=1.6pt, align=center,
    inner xsep=2.5pt, inner ysep=2pt, font=\srfont, text=srInk, minimum height=9.5mm},
  % client and edge (browser, CloudFront, S3)
  sr/client/.style={sr/base, fill=srBlue!10},
  % identity, scheduling, operations (Cognito, EventBridge, CloudWatch, SES)
  sr/service/.style={sr/base, fill=srPurple!8},
  % compute (Lambda functions, API Gateway)
  sr/compute/.style={sr/base, fill=srOrange!11},
  % queues (SQS)
  sr/queue/.style={sr/base, fill=srYellow!16},
  % data (DynamoDB)
  sr/data/.style={sr/base, fill=srGreen!11},
  % numbered step marker, placed on a path: node[sr/step]{3}
  sr/step/.style={circle, fill=srInk, text=white, font=\sffamily\bfseries\fontsize{6}{6}\selectfont,
    inner sep=0pt, minimum size=8.4pt},
  % arrows
  sr/flow/.style={draw=srInk, line width=0.55pt, -{Stealth[length=3.6pt,width=2.8pt]}},
  sr/async/.style={sr/flow, dash pattern=on 2.4pt off 1.4pt},
  sr/reply/.style={sr/flow, densely dotted},
  sr/aux/.style={draw=srMuted, line width=0.45pt, dotted, -{Stealth[length=3.2pt,width=2.6pt]}},
  % arrow and annotation labels
  sr/lbl/.style={font=\srfont, text=srInk, align=center, inner sep=1pt},
  sr/note/.style={font=\srfont, text=srMuted, align=center, inner sep=1pt},
}

% Then place a figure with
%   \begin{figure}[t]\centering% Fig. 2 (updated for the Final report): cloud architecture of SeatRelay.
% Source of truth: docs/ARCHITECTURE.md §1.2 (Figure A), §3–§7. Needs \input{figures/tikz-setup}.
% Size: about 15.8 cm x 6.5 cm (fits 0.95\textwidth of A4 with 2 cm margins; height <= 7 cm).
%
% Step numbers follow one write as in Figure 2 of the Preliminary report: (1) assets, (2) sign-in,
% (3) POST through the JWT authorizer to the api function (which records CMD#{id} = ACCEPTED),
% (4) enqueue in the workshop's FIFO group and answer 202, (5) the worker (single writer) runs
% decide(), (6) one TransactWriteItems guarded by the META version, (7) notifications via Streams,
% notifier and SES, (8a/8b) the per-offer timer re-enters the same FIFO group, (9) the browser
% polls GET /commands/{id}. Changed since the Preliminary report: the timer path 8a/8b replaces the
% expiry sweep (its old step 8), which now runs every 5 min only as a safety net (grey, unnumbered);
% notifications leave after the commit from the stream instead of from the worker. Nodes marked
% with an asterisk are new. Solid arrows: requests; dotted black: the 202 reply; dashed: queued
% messages and timers; grey: the safety net and the sweep's index query.
\begin{tikzpicture}[x=1cm, y=1cm]
  % ---------------------------------------------------------------- top band
  \node[sr/client,  minimum width=2.40cm] (cf)    at (1.20,6.00)  {\textbf{CloudFront + S3}\\private bucket, OAC};
  \node[sr/service, minimum width=1.90cm] (cog)   at (3.95,6.00)  {\textbf{Cognito}\\hosted UI};
  \node[sr/service, minimum width=1.90cm] (eb)    at (6.55,6.00)  {\textbf{EventBridge}\\Scheduler, 5 min};
  \node[sr/compute, minimum width=1.30cm] (sweep) at (8.75,6.00)  {\textbf{sweep}\\safety net};
  \node[sr/compute, minimum width=2.30cm] (fwd)   at (11.05,6.00) {\textbf{expiryForwarder}$^{\ast}$\\Lambda};
  \node[sr/queue,   minimum width=2.90cm] (delay) at (14.25,6.00) {\textbf{SQS expiry-delay}$^{\ast}$\\standard queue, one\\timer per offer, $\le$\,900\,s};
  % ---------------------------------------------------------------- request band
  \node[sr/client,  minimum width=1.65cm] (spa)   at (0.825,3.55)  {\textbf{Browser}\\SPA (React)};
  \node[sr/compute, minimum width=2.10cm] (gw)    at (5.82,3.55)  {\textbf{API Gateway}\\HTTP API\\JWT authorizer};
  \node[sr/compute, minimum width=1.30cm] (api)   at (8.00,3.55)  {\textbf{api}\\Lambda};
  \node[sr/queue,   minimum width=2.70cm] (fifo)  at (10.55,3.55) {\textbf{SQS FIFO}\\commands.fifo\\group = workshopId};
  \node[sr/compute, minimum width=2.90cm] (wk)    at (14.25,3.55) {\textbf{worker} Lambda\\tick-first decide()\\TransactWriteItems};
  % ---------------------------------------------------------------- data band
  \node[sr/data, minimum width=9.50cm, minimum height=1.10cm] (ddb) at (10.95,1.30)
    {\textbf{DynamoDB single table} (on-demand, PITR, TTL, Streams)\\
     CMD\#$\{$id$\}$: command results \enspace$\cdot$\enspace WS\#$\{$id$\}$: META with version lock,\\
     live L\# and history H\# items \enspace$\cdot$\enspace sparse DUE index (GSI1)};
  \node[sr/compute, minimum width=1.60cm] (nt)    at (3.80,1.30)  {\textbf{notifier}$^{\ast}$\\Lambda};
  \node[sr/service, minimum width=1.90cm] (ses)   at (0.95,1.30)  {\textbf{Amazon SES}\\e-mail};
  % ---------------------------------------------------------------- operations strip
  \node[sr/service, minimum width=15.70cm, minimum height=0.42cm, inner ysep=1pt, anchor=south west] (cw) at (0,0)
    {\textbf{CloudWatch}: EMF metrics, dashboard, alarms $\to$ SNS
     \enspace$\cdot$\enspace a DLQ behind every queue \enspace$\cdot$\enspace AWS Budgets (credits excluded)};

  % ---------------------------------------------------------------- (1) static assets, (2) sign-in
  \draw[sr/flow] (spa.north -| cf.south) -- node[sr/step, pos=0.45] {1} node[sr/lbl, left=5pt, pos=0.45] {assets} (cf.south);
  \draw[sr/flow, {Stealth[length=3.6pt,width=2.8pt]}-{Stealth[length=3.6pt,width=2.8pt]}]
    ($(spa.north east)+(-0.25,0)$) -- node[sr/step, pos=0.35] {2}
    node[sr/lbl, pos=0.5, right=6pt, align=left] {sign in,\\ID token} ($(cog.south west)+(0.40,0)$);
  % ---------------------------------------------------------------- (3) POST, 202, (9) poll
  \draw[sr/flow]  (spa.east |- 0,3.92) -- node[sr/step, pos=0.07] {3} node[sr/lbl, above, pos=0.58] {POST \ldots/commands} (gw.west |- 0,3.92);
  \draw[sr/reply] (gw.west |- 0,3.55) -- node[sr/lbl, above, pos=0.5] {202 Accepted} (spa.east |- 0,3.55);
  \draw[sr/flow]  (spa.east |- 0,3.18) -- node[sr/step, pos=0.07] {9} node[sr/lbl, above, pos=0.58] {GET /commands/$\{$id$\}$} (gw.west |- 0,3.18);
  \node[sr/note, anchor=north] at ($(spa.east |- 0,3.02)!0.5!(gw.west |- 0,3.02)$) {poll 300\,$\to$\,2000\,ms until\\APPLIED or REJECTED};
  % ---------------------------------------------------------------- (4) CMD record, (5) enqueue, (6) deliver
  \draw[sr/flow] (gw) -- (api);
  \draw[sr/flow] (api.south) -- node[sr/lbl, left=1pt, pos=0.32, align=right] {CMD\#$\{$id$\}$\\put, get} (api.south |- ddb.north);
  \draw[sr/flow] (api) -- node[sr/step] {4} (fifo);
  \draw[sr/flow] (fifo) -- node[sr/step] {5} (wk);
  % ---------------------------------------------------------------- worker reads and commits
  \coordinate (wkR) at ($(wk.south)+(0.85,0)$);
  \coordinate (wkL) at ($(wk.south)+(-0.85,0)$);
  \coordinate (wkT) at ($(wk.north)+(0.75,0)$);
  \draw[sr/flow] (wkR |- ddb.north) -- node[sr/lbl, left=1pt, pos=0.5, align=right] {Query L\#,\\consistent} (wkR);
  \draw[sr/flow] (wkL) -- node[sr/step, pos=0.32] {6} node[sr/lbl, left=3pt, pos=0.32, align=right] {one transaction,\\META version $= v$} (wkL |- ddb.north);
  % ---------------------------------------------------------------- (7) timer, (10) forward, (11) sweep
  \draw[sr/async] (wkT) -- node[sr/step, pos=0.45] {8a} node[sr/lbl, left=5pt, pos=0.45, align=right] {TICK per offer,\\before the commit} (wkT |- delay.south);
  \draw[sr/async] (delay) -- node[sr/step] {8b} (fwd);
  \coordinate (fwdB) at ($(fwd.south)+(-0.05,0)$);
  \draw[sr/async] (fwdB) -- node[sr/lbl, right=2pt, pos=0.45, align=left] {TICK,\\same group} (fwdB |- fifo.north);
  \draw[sr/flow, draw=srMuted] (eb) -- (sweep);
  \coordinate (swB) at ($(sweep.south)+(0.45,0)$);
  \draw[sr/async, draw=srMuted] (swB) -- node[sr/note, right=1pt, pos=0.45, align=left] {TICK if\\overdue} (swB |- fifo.north);
  \coordinate (gap) at ($(api.east)!0.5!(fifo.west)$);
  \draw[sr/aux] (gap |- sweep.south) -- node[sr/note, pos=0.22, left=1pt, align=right] {Query\\DUE index} (gap |- ddb.north);
  % ---------------------------------------------------------------- (12) notifications
  \draw[sr/flow] (ddb.west) -- node[sr/step] {7} node[sr/lbl, above=4pt] {Streams} (nt.east);
  \draw[sr/flow] (nt) -- node[sr/lbl, above] {send} (ses);
\end{tikzpicture}
\caption{...}\label{fig:arch}\end{figure}
%
% The figure files contain only a tikzpicture (no figure environment, no caption), so captions
% and labels stay in the report text. All diagram text is set at a fixed 7 pt sans size, so the
% figures have the same physical size whatever the document's base font size is. Widths are
% checked against a single-column A4 layout with 2 cm margins (\textwidth = 17 cm): every
% figure is at most 0.95\textwidth wide.
%
% Colours: the categorical slots of the validated palette used by the result charts
% (report/tools/make_figures.py), as light tints for node fills. Identity never rests on colour
% alone: every node names its service in text.

\usepackage{tikz}
\usetikzlibrary{arrows.meta,positioning,calc,fit,backgrounds,shapes.geometric,decorations.pathreplacing}

\definecolor{srBlue}{HTML}{2A78D6}
\definecolor{srOrange}{HTML}{EB6834}
\definecolor{srGreen}{HTML}{1BAF7A}
\definecolor{srYellow}{HTML}{EDA100}
\definecolor{srPurple}{HTML}{6250D6}
\definecolor{srInk}{HTML}{1F1F1D}
\definecolor{srMuted}{HTML}{52514E}
\definecolor{srGrid}{HTML}{D8D7D2}

\newcommand{\srfont}{\sffamily\fontsize{7}{8.2}\selectfont}
\newcommand{\srfontb}{\sffamily\bfseries\fontsize{7}{8.2}\selectfont}

\tikzset{
  sr/base/.style={draw=srInk!75, line width=0.5pt, rounded corners=1.6pt, align=center,
    inner xsep=2.5pt, inner ysep=2pt, font=\srfont, text=srInk, minimum height=9.5mm},
  % client and edge (browser, CloudFront, S3)
  sr/client/.style={sr/base, fill=srBlue!10},
  % identity, scheduling, operations (Cognito, EventBridge, CloudWatch, SES)
  sr/service/.style={sr/base, fill=srPurple!8},
  % compute (Lambda functions, API Gateway)
  sr/compute/.style={sr/base, fill=srOrange!11},
  % queues (SQS)
  sr/queue/.style={sr/base, fill=srYellow!16},
  % data (DynamoDB)
  sr/data/.style={sr/base, fill=srGreen!11},
  % numbered step marker, placed on a path: node[sr/step]{3}
  sr/step/.style={circle, fill=srInk, text=white, font=\sffamily\bfseries\fontsize{6}{6}\selectfont,
    inner sep=0pt, minimum size=8.4pt},
  % arrows
  sr/flow/.style={draw=srInk, line width=0.55pt, -{Stealth[length=3.6pt,width=2.8pt]}},
  sr/async/.style={sr/flow, dash pattern=on 2.4pt off 1.4pt},
  sr/reply/.style={sr/flow, densely dotted},
  sr/aux/.style={draw=srMuted, line width=0.45pt, dotted, -{Stealth[length=3.2pt,width=2.6pt]}},
  % arrow and annotation labels
  sr/lbl/.style={font=\srfont, text=srInk, align=center, inner sep=1pt},
  sr/note/.style={font=\srfont, text=srMuted, align=center, inner sep=1pt},
}
. Size: about 15.7 cm x 4.2 cm.
%
% Suggested caption:
%   User flow. A participant registers once with the acceptable sessions in preference order and
%   joins one queue per workshop. The first queued participant who accepts a session with a free
%   seat receives an offer that is held for holdSeconds; accepting confirms the seat, declining or
%   letting it expire drops that session from the preference list (R5). Every released seat is
%   offered to the next queued participant in the same decision (backfill). Organisers follow
%   seats, queue, backfills and distinct participants served on a live dashboard.
\begin{tikzpicture}[x=1cm, y=1cm]
  % ---------------------------------------------------------------- participant states
  \node[sr/client,  minimum width=2.35cm] (reg)  at (1.175,3.00) {\textbf{Register once}\\acceptable sessions,\\ranked};
  \node[sr/queue,   minimum width=1.90cm] (que)  at (4.60,3.00)  {\textbf{QUEUED}\\one queue per\\workshop};
  \node[sr/compute, minimum width=1.90cm] (off)  at (8.95,3.00)  {\textbf{OFFERED}\\seat held for\\holdSeconds};
  \node[sr/data,    minimum width=1.80cm] (con)  at (12.00,3.00) {\textbf{CONFIRMED}\\one seat per\\participant};
  \node[sr/service, minimum width=1.75cm, minimum height=0.60cm] (can) at (14.675,3.38) {\textbf{CANCELLED}};
  \node[sr/service, minimum width=1.75cm, minimum height=0.60cm] (end) at (14.675,2.45) {\textbf{ENDED}};

  \draw[sr/flow] (reg) -- node[sr/lbl, above=1pt] {register} (que);
  \draw[sr/flow] (que) -- node[sr/lbl, above=1pt, align=center] {free seat in an\\acceptable session} (off);
  \draw[sr/flow] (off) -- node[sr/lbl, above=1pt] {accept} (con);
  \draw[sr/flow] (con.east |- can.west) -- node[sr/lbl, above=1pt] {cancel} (can.west);
  \draw[sr/flow] (off.north) .. controls +(0,0.62) and +(0,0.55) .. node[sr/lbl, pos=0.5, above=0pt] {cancel (from any state)} (can.north);
  % decline or expiry: R5
  \draw[sr/flow] (off.south west) .. controls +(-0.4,-0.65) and +(0.6,-0.65) ..
    node[sr/lbl, below=1pt, align=center] {decline or expire:\\that session dropped (R5)} (que.south east);
  \draw[sr/flow] (off.south east) .. controls +(0.6,-0.65) and +(-1.4,0) ..
    node[sr/lbl, pos=0.5, below=1pt, align=center] {no acceptable session left} (end.west);

  % ---------------------------------------------------------------- backfill and organiser
  \node[sr/base, fill=srYellow!8, dash pattern=on 2pt off 1.2pt, minimum width=9.6cm, minimum height=0.5cm, anchor=west] (bf) at (3.65,0.95)
    {\textbf{Backfill}: a seat released by an expiry, a decline, a cancellation or a capacity increase\\
     is offered to the next queued participant who accepts that session, in the same decision};
  \coordinate (bfu) at ($(off.south)+(0.30,0)$);
  \draw[sr/async] (bfu |- bf.north) -- node[sr/lbl, right=1pt, pos=0.45] {offer} (bfu);
  \node[sr/service, minimum width=2.35cm] (org) at (1.175,0.95) {\textbf{Organiser}\\2--5 sessions,\\capacity, hold time};
  \draw[sr/flow] (org) -- node[sr/lbl, above=1pt] {capacity} (bf);
  \node[sr/data, text width=15.2cm, minimum height=0.5cm, anchor=north west, inner ysep=1.5pt] (dash) at (0,0.05)
    {\textbf{Live dashboard}: seats per session (capacity, confirmed, held, free) \enspace$\cdot$\enspace queue \enspace$\cdot$\enspace backfill log with latency
     \enspace$\cdot$\enspace distinct participants served \enspace$\cdot$\enspace attendance \enspace$\cdot$\enspace CSV export \enspace$\cdot$\enspace verify invariants};
  \draw[sr/flow] (org.south) -- (org.south |- dash.north);
\end{tikzpicture}

  \caption{User flow. A participant registers once with the acceptable sessions in preference order
  and joins one queue per workshop. The first queued participant who accepts a session with a free
  seat receives an offer held for \code{holdSeconds}; accepting confirms the seat, declining or
  letting it expire drops that session from the preference list (R5). Every released seat is offered
  to the next queued participant in the same decision (backfill).}
  \label{fig:flow}
\end{figure}

\paragraph{Evidence that repeated sessions occur.}
On the NUS Libraries calendar from 1~Aug to 16~Oct 2026, 11 of the 37 workshop events were
repeat sessions of two titles~\cite{nus-libcal}; on the SMU Libraries calendar in the same window,
two titles ran twice among 44 events~\cite{smu-libcal} (our counts on the public calendars). Most
titles on both calendars ran once, so calendars show that the situation exists but cannot estimate
how often an organiser faces it. That share, $p$, has to come from organiser interviews.

\paragraph{Market size.}
We size the market bottom-up (Table~\ref{tab:market}): the number of organisers who could run
multi-session sign-ups, times $p$, times the price $P$ of \sd120 per organiser-year
(\S\ref{sec:business}). Singapore has six autonomous universities and five
polytechnics~\cite{moe-au,moe-psei}; five of the universities list at least 693 student
organisations, a lower bound because counts given as \enquote{more than $X$} enter as $X$ and the
sixth university (SUSS) and the polytechnics are not counted. With $p = \text{\InterviewPEventWeightedPct}\,\%$ (event-weighted
share of organisers' events with two or more equivalent sessions;
\InterviewOrganisers{} organisers interviewed), the campus total addressable market is
$\text{TAM} = 693 \times p \times \text{\sd}120 = \text{\sd}83{,}160\,p$ per year, and
\sd167{,}160\,$p$ if the more than 700 CET training providers~\cite{ssg-credit-progress-2016} are
included. The serviceable market applies the share $s$ of organisers who are not bound to a closed
institutional system; the first-year obtainable market is NUS student organisations and pilot
units. Even at $p = 1$ the campus market is at most \sd83{,}160 a year, so the business case rests
on the time each organiser saves (\S\ref{sec:business}), not on market size.

\begin{table}[t]
\centering\small
\caption{Bottom-up market sizing (counts accessed 2026-10-17; details and verification levels in
\code{docs/market-sizing.md}).}\label{tab:market}
\begin{tabularx}{\textwidth}{@{}P{5.6cm}P{1.75cm}P{3.1cm}L@{}}
\toprule
Population & Count & Source & Estimation basis\\
\midrule
Student organisations, five autonomous universities (NUS 200, NTU 113, SMU 150, SUTD 80, SIT 150)
 & $\geq 693$ & \cite{nus-osa-getstarted,ntu-special-interest-clubs,ntu-su-council,ntu-clubs-societies,smu-osl-2022,sutd-fifth-row,sit-find-your-tribe}
 & official counts as lower bounds; NTU counted by us from its directory pages\\
CET training providers & $> 700$ (2016) & \cite{ssg-credit-progress-2016}
 & latest official count found; reported separately\\
Learners in SkillsFuture-supported training, 2025 & 606{,}000 & \cite{ssg-10th-anniversary-2026}
 & demand context only\\
Organisers with multi-session sign-ups & $693 \times p$ & organiser interviews
 & $p$: \InterviewPEventWeightedPct\\
Price $P$ per organiser & \sd120/year & \S\ref{sec:business} & \sd10 per month\\
\bottomrule
\end{tabularx}
\end{table}

\paragraph{Competitors.}
Table~\ref{tab:competitors} summarises the vendors' documentation. Eventbrite relays a released
ticket automatically with a claim window, but its inventory and waitlist are per date or ticket
type; no tool we scanned ties several acceptable sessions to one queue position with one seat per
person, which is what removes duplicate holds.

\begin{table}[t]
\centering\footnotesize
\caption{Competitor scan from vendor documentation read on 2026-10-17 (not our tests). n/d: not
documented in the pages we could read.}\label{tab:competitors}
\begin{tabularx}{\textwidth}{@{}P{2.45cm}P{3.0cm}P{2.2cm}P{2.6cm}P{2.3cm}L@{}}
\toprule
Tool & Several sessions & Waitlist & Timed relay & One seat across sessions & Organiser price\\
\midrule
Eventbrite~\cite{eventbrite-multidate,eventbrite-waitlist-feature,eventbrite-pricing-sg}
 & recurring event; inventory per date & per ticket type or date & yes, with claim window & no (per date) & free events free; fees on paid tickets\\
Luma~\cite{luma-multisession,luma-waitlist,luma-pricing}
 & workaround: cloned events & yes & no (host approves) & no & free; 5\,\% on paid events\\
SignUpGenius~\cite{signupgenius-waitlist,signupgenius-pricing}
 & one slot per session & per slot (paid) & relay without a hold & n/d & \usd22.49/month (Gold, annual)\\
Peatix~\cite{peatix-recurring,peatix-waitlist,peatix-pricing-sg}
 & up to 50 dates & no automated waitlist & no & no & free events free; 4.9\,\% + \sd0.99/ticket\\
Calendly group~\cite{calendly-group,calendly-pricing}
 & slots with invitee limit & n/d & no & n/d & paid plans; \usd10/seat/month (annual)\\
Google / MS Forms~\cite{google-forms-responses,ms-forms-settings}
 & one form per session & no & no & manual / per form & pricing not checked\\
\textbf{SeatRelay} & 2--5 sessions, one queue & yes & yes, timed hold & yes (C2) & free tier; \sd10/month\\
\bottomrule
\end{tabularx}
\end{table}
